% Generated draft: requires mathematical and source review.
\documentclass[11pt]{article}
\usepackage[margin=25mm]{geometry}
\usepackage{fontspec}
\setmainfont{Latin Modern Roman}
\usepackage{amsmath,amssymb,mathtools,unicode-math}
\setmathfont{Latin Modern Math}
\usepackage{enumitem,needspace,xurl,hyperref}
\setlength{\parindent}{0pt}
\setlength{\parskip}{5pt}
\setlength{\emergencystretch}{3em}
\newcommand{\N}{\mathbb{N}}
\newcommand{\Z}{\mathbb{Z}}
\newcommand{\Q}{\mathbb{Q}}
\newcommand{\R}{\mathbb{R}}
\newcommand{\C}{\mathbb{C}}
\newcommand{\F}{\mathbb{F}}
\newcommand{\A}{\mathbb{A}}
\newcommand{\PP}{\mathbb{P}}
\newcommand{\E}{\mathbb{E}}
\newcommand{\eps}{\varepsilon}
\newcommand{\dd}{\,\mathrm{d}}
\begin{document}
\section*{UnsolvedMath ID: 20001283}
\textbf{Source problem code:} catalogue entry
\subsection{Definitions and mathematical statement}
Let $n\geq3$ and $2\leq k\leq n-1$, and let $G(n,k)$ denote the $k$-dimensional linear subspaces of $\R^n$. For sets in a common affine space, a positive homothety has the form $x\mapsto\lambda x+a$, $\lambda>0$; a translation has $\lambda=1$. The source does not explicitly fix the sign convention for homothety.

The quoted theorem of Rogers (1965) concerns compact convex sets $K,L\subset\R^n$ with $0\in\operatorname{relint}K\cap\operatorname{relint}L$: if, for every $S\in G(n,k)$, the corresponding sections $K\cap S$ and $L\cap S$ are homothetic, then $K$ and $L$ are homothetic; if every corresponding pair is translates, then $K$ and $L$ are translates. Relative interior is interior within the affine hull, so the quoted convex theorem is not silently restricted to full-dimensional sets.

R. J. Gardner asks whether these two assertions continue to hold for star bodies, or for a still broader class. Under the usual convention, a star body about $0$ is a compact set with continuous positive radial function $\rho_K:S^{n-1}\to(0,\infty)$ and $K\cap\R_{\geq0}u=[0,\rho_K(u)]u$. For such bodies, the homothety question is whether
\[
 \bigl[\forall S\in G(n,k)\ \exists\lambda_S>0,\ a_S\in S:\
 K\cap S=a_S+\lambda_S(L\cap S)\bigr]
 \ \Longrightarrow\ \exists\lambda>0,\ a\in\R^n:\ K=a+\lambda L.
\]
The translation question replaces every $\lambda_S$ and the concluding $\lambda$ by $1$. The section translations and homothety centers may depend on $S$; none is assumed centered at the origin. The unspecified ``more generally still'' is a request for possible extensions, not a defined extra class.

The supplied machine aid reports a centered-dilation partial result. For star bodies set
\[
 f(u)=\log\rho_K(u)-\log\rho_L(u),\quad
 D(K,L)=\inf_{c\in\R}\sup_{u\in S^{n-1}}|f(u)-c|,
\]
\[
 D_S(K,L)=\inf_{c\in\R}\sup_{u\in S\cap S^{n-1}}|f(u)-c|.
\]
It states that for $2\leq k\leq n$,
\[
 D(K,L)=\max_{S\in G(n,k)}D_S(K,L).
\]
Thus, for $\varepsilon\geq0$, sectionwise error at most $\varepsilon$ after optimally chosen centered dilations yields a single global centered dilation with the same log-radial error: some $c\in\R$ satisfies $e^{-\varepsilon}\leq\rho_K(u)/(e^c\rho_L(u))\leq e^{\varepsilon}$ for all $u$. At $\varepsilon=0$ this concerns centered section dilations, not arbitrary translated section homotheties.

The same aid reports that, if $K=-K$ and $L=-L$, uniqueness of the symmetry center forces section homotheties to fix $0$, so the homothety branch holds for this symmetric subclass. Sectionwise translations of such symmetric bodies force $K=L$; the aid says this translation conclusion also holds for $k=1$. These are attributed restricted statements and do not settle the original nonsymmetric extension question.
\subsection{Short English statement}
Rogers' convex theorem says that matching corresponding central sections by homotheties, or by translations, forces the same relation between the whole bodies. Gardner asks whether convexity can be replaced by the star-body condition. The supplied partial result glues dilations centered at the origin with no loss of radial error, and handles symmetric bodies; arbitrary section translations for nonsymmetric star bodies remain outside that result.
\subsection{Sources}
\begin{itemize}
\item Local UnsolvedMath registry, permanent ID 20001283, exact source code catalogue entry; canonical \texttt{aim-convex-geometry-notes.json}, \texttt{notes[14]}, source item 9, tag ``problem'', workshop \emph{Mahler's conjecture and duality in convex geometry}. Supplied URL: \url{https://aimath.org/WWN/mahlerduality/mahlerduality.pdf}.
\item Original attribution: R. J. Gardner; quoted theorem: Rogers, 1965. Original citations: [G, Problem 7.1] and [G, Note 7.1]. No expansion of [G] or full Rogers bibliographic entry is supplied. The verification identifies the official 2010 AIM workshop list and announces three OCR corrections, but truncates before listing them.
\item Supplied research aid: PARTIAL-PROGRESS; registry status ``partially\_solved'', statement status ``exact''; attempt 1, result ``partial\_result'', problem status at run ``open'', recorded rigor ``proved'', novelty confidence ``low'', candidate type ``stability\_identity''. The aid is explicitly machine generated, not peer reviewed, and requires independent verification before citation. Original remarks and literature field are empty. No external source was visited.
\end{itemize}
\end{document}
